% Options for packages loaded elsewhere
\PassOptionsToPackage{unicode}{hyperref}
\PassOptionsToPackage{hyphens}{url}
%
\documentclass[
  12pt,
  a4paper,
  DIV=11,
  numbers=noendperiod]{scrreprt}

\usepackage{amsmath,amssymb}
\usepackage{setspace}
\usepackage{iftex}
\ifPDFTeX
  \usepackage[T1]{fontenc}
  \usepackage[utf8]{inputenc}
  \usepackage{textcomp} % provide euro and other symbols
\else % if luatex or xetex
  \usepackage{unicode-math}
  \defaultfontfeatures{Scale=MatchLowercase}
  \defaultfontfeatures[\rmfamily]{Ligatures=TeX,Scale=1}
\fi
\usepackage{lmodern}
\ifPDFTeX\else  
    % xetex/luatex font selection
\fi
% Use upquote if available, for straight quotes in verbatim environments
\IfFileExists{upquote.sty}{\usepackage{upquote}}{}
\IfFileExists{microtype.sty}{% use microtype if available
  \usepackage[]{microtype}
  \UseMicrotypeSet[protrusion]{basicmath} % disable protrusion for tt fonts
}{}
\usepackage{xcolor}
\setlength{\emergencystretch}{3em} % prevent overfull lines
\setcounter{secnumdepth}{5}
% Make \paragraph and \subparagraph free-standing
\makeatletter
\ifx\paragraph\undefined\else
  \let\oldparagraph\paragraph
  \renewcommand{\paragraph}{
    \@ifstar
      \xxxParagraphStar
      \xxxParagraphNoStar
  }
  \newcommand{\xxxParagraphStar}[1]{\oldparagraph*{#1}\mbox{}}
  \newcommand{\xxxParagraphNoStar}[1]{\oldparagraph{#1}\mbox{}}
\fi
\ifx\subparagraph\undefined\else
  \let\oldsubparagraph\subparagraph
  \renewcommand{\subparagraph}{
    \@ifstar
      \xxxSubParagraphStar
      \xxxSubParagraphNoStar
  }
  \newcommand{\xxxSubParagraphStar}[1]{\oldsubparagraph*{#1}\mbox{}}
  \newcommand{\xxxSubParagraphNoStar}[1]{\oldsubparagraph{#1}\mbox{}}
\fi
\makeatother


\providecommand{\tightlist}{%
  \setlength{\itemsep}{0pt}\setlength{\parskip}{0pt}}\usepackage{longtable,booktabs,array}
\usepackage{calc} % for calculating minipage widths
% Correct order of tables after \paragraph or \subparagraph
\usepackage{etoolbox}
\makeatletter
\patchcmd\longtable{\par}{\if@noskipsec\mbox{}\fi\par}{}{}
\makeatother
% Allow footnotes in longtable head/foot
\IfFileExists{footnotehyper.sty}{\usepackage{footnotehyper}}{\usepackage{footnote}}
\makesavenoteenv{longtable}
\usepackage{graphicx}
\makeatletter
\newsavebox\pandoc@box
\newcommand*\pandocbounded[1]{% scales image to fit in text height/width
  \sbox\pandoc@box{#1}%
  \Gscale@div\@tempa{\textheight}{\dimexpr\ht\pandoc@box+\dp\pandoc@box\relax}%
  \Gscale@div\@tempb{\linewidth}{\wd\pandoc@box}%
  \ifdim\@tempb\p@<\@tempa\p@\let\@tempa\@tempb\fi% select the smaller of both
  \ifdim\@tempa\p@<\p@\scalebox{\@tempa}{\usebox\pandoc@box}%
  \else\usebox{\pandoc@box}%
  \fi%
}
% Set default figure placement to htbp
\def\fps@figure{htbp}
\makeatother

\usepackage{indentfirst}
\usepackage{float}
\floatplacement{figure}{H}
\IfFileExists{plex-otf.sty}{
  %% Full TeXlive
  \usepackage[
    % math,
    RM={Scale=0.94},SS={Scale=0.94},SScon={Scale=0.94},TT={Scale=MatchLowercase,FakeStretch=0.9},
    DefaultFeatures={Ligatures=Common}
  ]{plex-otf}
  % \usepackage[Scale=MatchUppercase]{juliamono}
}{
  %% TinyTeX
  \usepackage{libertine}
}
\KOMAoption{captions}{tableheading}
\makeatletter
\@ifpackageloaded{caption}{}{\usepackage{caption}}
\AtBeginDocument{%
\ifdefined\contentsname
  \renewcommand*\contentsname{Содержание}
\else
  \newcommand\contentsname{Содержание}
\fi
\ifdefined\listfigurename
  \renewcommand*\listfigurename{Список иллюстраций}
\else
  \newcommand\listfigurename{Список иллюстраций}
\fi
\ifdefined\listtablename
  \renewcommand*\listtablename{Список таблиц}
\else
  \newcommand\listtablename{Список таблиц}
\fi
\ifdefined\figurename
  \renewcommand*\figurename{Рисунок}
\else
  \newcommand\figurename{Рисунок}
\fi
\ifdefined\tablename
  \renewcommand*\tablename{Таблица}
\else
  \newcommand\tablename{Таблица}
\fi
}
\@ifpackageloaded{float}{}{\usepackage{float}}
\floatstyle{ruled}
\@ifundefined{c@chapter}{\newfloat{codelisting}{h}{lop}}{\newfloat{codelisting}{h}{lop}[chapter]}
\floatname{codelisting}{Список}
\newcommand*\listoflistings{\listof{codelisting}{Листинги}}
\makeatother
\makeatletter
\makeatother
\makeatletter
\@ifpackageloaded{caption}{}{\usepackage{caption}}
\@ifpackageloaded{subcaption}{}{\usepackage{subcaption}}
\makeatother

\ifLuaTeX
\usepackage[bidi=basic]{babel}
\else
\usepackage[bidi=default]{babel}
\fi
\babelprovide[main,import]{russian}
\babelprovide[import]{english}
% get rid of language-specific shorthands (see #6817):
\let\LanguageShortHands\languageshorthands
\def\languageshorthands#1{}
\usepackage[style=gost-numeric,backend=biber,langhook=extras,autolang=other*]{biblatex}
\addbibresource{bib/cite.bib}
\usepackage{csquotes}
\usepackage{bookmark}

\IfFileExists{xurl.sty}{\usepackage{xurl}}{} % add URL line breaks if available
\urlstyle{same} % disable monospaced font for URLs
\hypersetup{
  pdftitle={Отчёт по лабораторной работе №10},
  pdfauthor={Пакавира Арсениу Висенте Луиш},
  pdflang={ru-RU},
  hidelinks,
  pdfcreator={LaTeX via pandoc}}


\title{Отчёт по лабораторной работе №10}
\usepackage{etoolbox}
\makeatletter
\providecommand{\subtitle}[1]{% add subtitle to \maketitle
  \apptocmd{\@title}{\par {\large #1 \par}}{}{}
}
\makeatother
\subtitle{Дисциплина: Администрирование сетевых подсистем}
\author{Пакавира Арсениу Висенте Луиш}
\date{}

\begin{document}
\maketitle

\renewcommand*\contentsname{Содержание}
{
\setcounter{tocdepth}{1}
\tableofcontents
}
\listoffigures
\listoftables

\setstretch{1.5}
\chapter{Цель
работы}\label{ux446ux435ux43bux44c-ux440ux430ux431ux43eux442ux44b}

Приобретение практических навыков по конфигурированию SMTP-сервера в
части настройки аутентификации.

\chapter{Выполнение лабораторной
работы}\label{ux432ux44bux43fux43eux43bux43dux435ux43dux438ux435-ux43bux430ux431ux43eux440ux430ux442ux43eux440ux43dux43eux439-ux440ux430ux431ux43eux442ux44b}

\section{Настройка LMTP в
Dovecote}\label{ux43dux430ux441ux442ux440ux43eux439ux43aux430-lmtp-ux432-dovecote}

Вход в систему выполнен под пользователем с последующим переходом в
режим суперпользователя с помощью команды \texttt{sudo\ -i}, что
подтверждается приглашением \texttt{root@server} (\textbf{?@fig-1}). Это
обеспечивает права для изменения конфигурации почтовых служб.

\includegraphics[width=0.7\linewidth,height=\textheight,keepaspectratio]{image/1.png}
\href{https://github.com/ArsenioPacavira99/study_2025-2026_net-os-admin-12/tree/master/labs/lab10/report/image/1.png}{рис.1}

В дополнительном терминале запущен мониторинг журнала почтовой службы
командой \texttt{tail\ -f\ /var/log/maillog}, что позволяет в реальном
времени отслеживать события Postfix и Dovecot, включая запуск демонов и
обработку соединений (\textbf{?@fig-1}).

В конфигурационном файле \texttt{/etc/dovecot/dovecot.conf} в список
обслуживаемых протоколов добавлен протокол LMTP посредством строки
\texttt{protocols\ =\ imap\ pop3\ lmtp}, что расширяет функциональность
Dovecot для работы с LMTP (\textbf{?@fig-2}).

\includegraphics[width=0.7\linewidth,height=\textheight,keepaspectratio]{image/2.png}
\href{https://github.com/ArsenioPacavira99/study_2025-2026_net-os-admin-12/tree/master/labs/lab10/report/image/2.png}{рис.2}

В файле \texttt{/etc/dovecot/conf.d/10-master.conf} выполнена настройка
сервиса \texttt{lmtp}, где определён unix-сокет
\texttt{/var/spool/postfix/private/dovecot-lmtp} с параметрами доступа
\texttt{user\ =\ postfix}, \texttt{group\ =\ postfix},
\texttt{mode\ =\ 0600}, обеспечивающими взаимодействие Dovecot с Postfix
через локальный сокет (\textbf{?@fig-3}).

\includegraphics[width=0.7\linewidth,height=\textheight,keepaspectratio]{image/3.png}
\href{https://github.com/ArsenioPacavira99/study_2025-2026_net-os-admin-12/tree/master/labs/lab10/report/image/3.png}{рис.3}

В конфигурации Postfix параметр \texttt{mailbox\_transport}
переопределён командой
\texttt{postconf\ -e\ \textquotesingle{}mailbox\_transport\ =\ lmtp:unix:private/dovecot-lmtp\textquotesingle{}},
что направляет доставку почты через unix-сокет Dovecot вместо прямой
локальной доставки (\textbf{?@fig-4}).

\includegraphics[width=0.7\linewidth,height=\textheight,keepaspectratio]{image/4.png}
\href{https://github.com/ArsenioPacavira99/study_2025-2026_net-os-admin-12/tree/master/labs/lab10/report/image/4.png}{рис.41}

В файле \texttt{/etc/dovecot/conf.d/10-auth.conf} установлен формат
имени пользователя \texttt{auth\_username\_format\ =\ \%Ln}, что
обеспечивает аутентификацию по логину без доменной части
(рис.~\ref{fig-5}).

\begin{figure}

\centering{

\includegraphics[width=0.7\linewidth,height=\textheight,keepaspectratio]{image/5.png}

}

\caption{\label{fig-5}Настройка формата имени пользователя в Dovecot}

\end{figure}%

После внесения изменений выполнен перезапуск служб \texttt{postfix} и
\texttt{dovecot} командами \texttt{systemctl\ restart\ postfix} и
\texttt{systemctl\ restart\ dovecot}, что активирует обновлённые
параметры конфигурации (рис.~\ref{fig-6}).

\begin{figure}

\centering{

\includegraphics[width=0.7\linewidth,height=\textheight,keepaspectratio]{image/6.png}

}

\caption{\label{fig-6}Перезапуск служб Postfix и Dovecot}

\end{figure}%

С клиента отправлено тестовое письмо командой
\texttt{echo\ .\textbar{}\ mail\ -s\ "LMTP\ test"\ bahis@bahis.net}, что
инициирует передачу сообщения через SMTP на сервер (рис.~\ref{fig-8}).

\begin{figure}

\centering{

\includegraphics[width=0.7\linewidth,height=\textheight,keepaspectratio]{image/8.png}

}

\caption{\label{fig-8}Отправка тестового письма с клиента}

\end{figure}%

На сервере выполнена проверка почтового ящика пользователя с
использованием переменной \texttt{MAIL=\textasciitilde{}/Maildir/} и
команды \texttt{mail}. В списке сообщений отображается письмо с темой
«LMTP test», что подтверждает корректную доставку через LMTP в каталог
Maildir (рис.~\ref{fig-9}).

\begin{figure}

\centering{

\includegraphics[width=0.7\linewidth,height=\textheight,keepaspectratio]{image/9.png}

}

\caption{\label{fig-9}Просмотр содержимого Maildir и подтверждение
доставки письма}

\end{figure}%

\section{Настройка
SMTP-аутентификации}\label{ux43dux430ux441ux442ux440ux43eux439ux43aux430-smtp-ux430ux443ux442ux435ux43dux442ux438ux444ux438ux43aux430ux446ux438ux438}

В файле \texttt{/etc/dovecot/conf.d/10-master.conf} определена служба
аутентификации \texttt{service\ auth}, в которой задан unix-сокет
\texttt{/var/spool/postfix/private/auth} с параметрами
\texttt{user\ =\ postfix}, \texttt{group\ =\ postfix},
\texttt{mode\ =\ 0660}. Это обеспечивает доступ Postfix к механизму
SASL-аутентификации Dovecot через локальный сокет. Дополнительно
определён \texttt{unix\_listener\ auth-userdb} с правами
\texttt{mode\ =\ 0600} и пользователем \texttt{dovecot}, что
ограничивает доступ к базе пользователей процессом Dovecot
(рис.~\ref{fig-10}).

\begin{figure}

\centering{

\includegraphics[width=0.7\linewidth,height=\textheight,keepaspectratio]{image/10.png}

}

\caption{\label{fig-10}Определение службы аутентификации service auth в
Dovecot}

\end{figure}%

В Postfix задан тип SASL-аутентификации \texttt{dovecot} и путь к
unix-сокету \texttt{private/auth} с помощью параметров
\texttt{smtpd\_sasl\_type} и \texttt{smtpd\_sasl\_path}. Это настраивает
smtpd на использование Dovecot как backend-аутентификатора
(рис.~\ref{fig-11}).

\begin{figure}

\centering{

\includegraphics[width=0.7\linewidth,height=\textheight,keepaspectratio]{image/11.png}

}

\caption{\label{fig-11}Настройка smtpd\_sasl\_type и smtpd\_sasl\_path в
Postfix}

\end{figure}%

Параметр \texttt{smtpd\_recipient\_restrictions} настроен с
последовательностью: \texttt{reject\_unknown\_recipient\_domain} ---
отклонение писем для неизвестных доменов; \texttt{permit\_mynetworks}
--- разрешение для доверенной сети;
\texttt{reject\_non\_fqdn\_recipient} --- запрет некорректных адресов;
\texttt{reject\_unauth\_destination} --- запрет релея для неразрешённых
направлений; \texttt{reject\_unverified\_recipient} --- проверка
существования получателя; \texttt{permit} --- окончательное разрешение
при выполнении условий. Это предотвращает использование сервера как SMTP
relay (рис.~\ref{fig-11}).

В параметре \texttt{mynetworks} установлено значение
\texttt{127.0.0.0/8}, что ограничивает приём почты только локальными
соединениями (рис.~\ref{fig-11}).

В файле \texttt{/etc/postfix/master.cf} для сервиса \texttt{smtp}
добавлены параметры \texttt{-o\ smtpd\_sasl\_auth\_enable=yes} и
модифицированы ограничения получателей, что временно включает
аутентификацию на порту 25 для тестирования (рис.~\ref{fig-12}).

\begin{figure}

\centering{

\includegraphics[width=0.7\linewidth,height=\textheight,keepaspectratio]{image/12.png}

}

\caption{\label{fig-12}Модификация сервиса smtp в master.cf для
включения SASL}

\end{figure}%

После внесения изменений выполнен перезапуск служб \texttt{postfix} и
\texttt{dovecot}, что активирует новую конфигурацию аутентификации
(рис.~\ref{fig-13}).

\begin{figure}

\centering{

\includegraphics[width=0.7\linewidth,height=\textheight,keepaspectratio]{image/13.png}

}

\caption{\label{fig-13}Перезапуск служб Postfix и Dovecot}

\end{figure}%

На клиенте установлен пакет \texttt{telnet} с использованием
\texttt{dnf}, что позволяет выполнить ручное тестирование
SMTP-соединения (рис.~\ref{fig-14}).

\begin{figure}

\centering{

\includegraphics[width=0.7\linewidth,height=\textheight,keepaspectratio]{image/14.png}

}

\caption{\label{fig-14}Установка telnet на клиенте}

\end{figure}%

Сформирована строка аутентификации в формате base64 командой
\texttt{printf\ \textquotesingle{}bahis\textbackslash{}x00bahis\textbackslash{}x00123456\textquotesingle{}\ \textbar{}\ base64},
затем выполнено подключение к SMTP-серверу по порту 25. После команды
\texttt{EHLO\ test} сервер объявил поддержку \texttt{AUTH\ PLAIN}.
Команда \texttt{AUTH\ PLAIN\ \textless{}base64\textgreater{}}
завершилась ответом \texttt{235\ 2.7.0\ Authentication\ successful}, что
подтверждает корректную работу SASL-аутентификации через Dovecot
(рис.~\ref{fig-15}).

\begin{figure}

\centering{

\includegraphics[width=0.7\linewidth,height=\textheight,keepaspectratio]{image/15.png}

}

\caption{\label{fig-15}Тестирование SMTP-аутентификации через telnet}

\end{figure}%

\section{Настройка SMTP over
TLS}\label{ux43dux430ux441ux442ux440ux43eux439ux43aux430-smtp-over-tls}

На сервере выполнено копирование временного сертификата и закрытого
ключа Dovecot из каталога \texttt{/etc/pki/dovecot} в соответствующие
подкаталоги \texttt{/etc/pki/tls/certs} и \texttt{/etc/pki/tls/private},
после чего в Postfix заданы параметры \texttt{smtpd\_tls\_cert\_file},
\texttt{smtpd\_tls\_key\_file},
\texttt{smtpd\_tls\_session\_cache\_database}, а также уровни
безопасности \texttt{smtpd\_tls\_security\_level\ =\ may} и
\texttt{smtp\_tls\_security\_level\ =\ may}, что активирует поддержку
TLS для входящих и исходящих SMTP-соединений (рис.~\ref{fig-16}).

\begin{figure}

\centering{

\includegraphics[width=0.7\linewidth,height=\textheight,keepaspectratio]{image/16.png}

}

\caption{\label{fig-16}Настройка TLS-параметров в Postfix}

\end{figure}%

В файле \texttt{/etc/postfix/master.cf} оставлен стандартный сервис
\texttt{smtp}, а также добавлен сервис \texttt{submission} на порту 587
с параметрами \texttt{-o\ smtpd\_tls\_security\_level=encrypt},
\texttt{-o\ smtpd\_sasl\_auth\_enable=yes} и ограничениями получателей,
что обеспечивает обязательное использование TLS и аутентификации для
клиентов (рис.~\ref{fig-18}).

\begin{figure}

\centering{

\includegraphics[width=0.7\linewidth,height=\textheight,keepaspectratio]{image/18.png}

}

\caption{\label{fig-18}Добавление сервиса submission в master.cf}

\end{figure}%

В межсетевом экране разрешена служба \texttt{smtp-submission} командами
\texttt{firewall-cmd\ -\/-add-service=smtp-submission} и с сохранением
правила \texttt{-\/-permanent}, после чего выполнена перезагрузка
конфигурации брандмауэра и перезапуск Postfix (рис.~\ref{fig-19}).

\begin{figure}

\centering{

\includegraphics[width=0.7\linewidth,height=\textheight,keepaspectratio]{image/19.png}

}

\caption{\label{fig-19}Настройка firewall и перезапуск Postfix}

\end{figure}%

С клиента выполнено подключение к серверу по порту 587 с использованием
\texttt{openssl\ s\_client\ -starttls\ smtp\ -crlf\ -connect\ server.bahis.net:587}.
В процессе установлено TLS-соединение, отображена информация о
самоподписанном сертификате, после чего сервер успешно принял команды
\texttt{EHLO} и объявил поддержку \texttt{AUTH\ PLAIN}
(рис.~\ref{fig-20}).

\begin{figure}

\centering{

\includegraphics[width=0.7\linewidth,height=\textheight,keepaspectratio]{image/20.png}

}

\caption{\label{fig-20}Установление SMTP over TLS соединения через
openssl}

\end{figure}%

Проверка аутентификации выполнена командой
\texttt{AUTH\ PLAIN\ \textless{}base64\textgreater{}}, сервер вернул
ответ \texttt{235\ 2.7.0\ Authentication\ successful}, что подтверждает
корректную работу SMTP over TLS с SASL-аутентификацией
(рис.~\ref{fig-21}).

\begin{figure}

\centering{

\includegraphics[width=0.7\linewidth,height=\textheight,keepaspectratio]{image/21.png}

}

\caption{\label{fig-21}Успешная SMTP-аутентификация через порт 587}

\end{figure}%

В почтовом клиенте Evolution настроен SMTP-сервер с использованием порта
587, режима STARTTLS и обычного пароля, после чего отправлено тестовое
сообщение с темой «test 5» (рис.~\ref{fig-22}).

\begin{figure}

\centering{

\includegraphics[width=0.7\linewidth,height=\textheight,keepaspectratio]{image/22.png}

}

\caption{\label{fig-22}Отправка письма через Evolution по SMTP over TLS}

\end{figure}%

На сервере выполнена проверка почтового ящика командой
\texttt{MAIL=\textasciitilde{}/Maildir/\ mail}, где отображается
сообщение с темой «test 5», что подтверждает корректную доставку письма
через защищённое соединение SMTP over TLS (рис.~\ref{fig-23}).

\begin{figure}

\centering{

\includegraphics[width=0.7\linewidth,height=\textheight,keepaspectratio]{image/23.png}

}

\caption{\label{fig-23}Подтверждение доставки письма в Maildir}

\end{figure}%

\section{Внесение изменений в настройки
внутреннего}\label{ux432ux43dux435ux441ux435ux43dux438ux435-ux438ux437ux43cux435ux43dux435ux43dux438ux439-ux432-ux43dux430ux441ux442ux440ux43eux439ux43aux438-ux432ux43dux443ux442ux440ux435ux43dux43dux435ux433ux43e}

На виртуальной машине \textbf{server} выполнен переход в каталог
\texttt{/vagrant/provision/server}, после чего конфигурационные файлы
\texttt{dovecot.conf}, \texttt{10-master.conf}, \texttt{10-auth.conf}, а
также \texttt{master.cf} Postfix скопированы в соответствующие
подкаталоги \texttt{mail/etc/dovecot/} и \texttt{mail/etc/postfix/}. Это
обеспечивает сохранение рабочей конфигурации служб для последующего
автоматического развёртывания через provisioning (рис.~\ref{fig-24}).

\begin{figure}

\centering{

\includegraphics[width=0.7\linewidth,height=\textheight,keepaspectratio]{image/24.png}

}

\caption{\label{fig-24}Копирование конфигурационных файлов Dovecot и
Postfix в каталог provision}

\end{figure}%

В файл \texttt{/vagrant/provision/server/mail.sh} внесены изменения
расширенной конфигурации SMTP-сервера: добавлена установка пакетов
\texttt{postfix}, \texttt{dovecot}, \texttt{telnet}; копирование
конфигурации в \texttt{/etc}; настройка firewall для служб
\texttt{smtp}, \texttt{pop3}, \texttt{pop3s}, \texttt{imap},
\texttt{imaps}, \texttt{smtp-submission}; параметры Postfix
(\texttt{mydomain}, \texttt{myorigin}, \texttt{inet\_protocols},
\texttt{inet\_interfaces}, \texttt{mydestination}); поддержка Maildir
(\texttt{home\_mailbox\ =\ Maildir/}); SASL-аутентификация через Dovecot
(\texttt{smtpd\_sasl\_type}, \texttt{smtpd\_sasl\_path}); ограничения
\texttt{smtpd\_recipient\_restrictions}; параметр \texttt{mynetworks}; а
также конфигурация SMTP over TLS с указанием сертификата, ключа и уровня
безопасности. В конце скрипта выполняются
\texttt{postfix\ set-permissions}, \texttt{restorecon}, перезапуск служб
Postfix и Dovecot (рис.~\ref{fig-25}).

\begin{figure}

\centering{

\includegraphics[width=0.7\linewidth,height=\textheight,keepaspectratio]{image/25.png}

}

\caption{\label{fig-25}Расширенная конфигурация SMTP-сервера в mail.sh
(server)}

\end{figure}%

В файл \texttt{/vagrant/provision/client/mail.sh} добавлена установка
пакета \texttt{telnet} командой \texttt{dnf\ -y\ install\ telnet}, что
обеспечивает возможность тестирования SMTP-соединения с клиента в
процессе автоматической настройки виртуальной машины
(рис.~\ref{fig-26}).

\begin{figure}

\centering{

\includegraphics[width=0.7\linewidth,height=\textheight,keepaspectratio]{image/26.png}

}

\caption{\label{fig-26}Добавление установки telnet в mail.sh (client)}

\end{figure}%

\chapter{Выводы}\label{ux432ux44bux432ux43eux434ux44b}

В ходе работы выполнена комплексная настройка почтового сервера на базе
Postfix и Dovecot с поддержкой LMTP, SASL-аутентификации и SMTP over
TLS.

Реализована доставка почты через LMTP-сокет Dovecot, что обеспечило
корректную интеграцию MTA и MDA. Настроена SASL-аутентификация через
Dovecot, исключающая несанкционированное использование сервера в
качестве SMTP-relay. Ограничения \texttt{smtpd\_recipient\_restrictions}
и параметр \texttt{mynetworks} обеспечили контроль приёма сообщений.

Настроена поддержка TLS с использованием сертификата Dovecot и
реализован сервис submission (порт 587) с обязательным шифрованием и
аутентификацией. Проверка через \texttt{openssl} и \texttt{AUTH\ PLAIN}
подтвердила успешное установление защищённого соединения.

Отправка сообщений через telnet и почтовый клиент Evolution показала
корректную работу SMTP over TLS и доставку писем в Maildir пользователя.

Конфигурация интегрирована в механизм provisioning виртуальной машины,
что обеспечивает воспроизводимость настройки внутреннего окружения.

\chapter{Контрольные
вопросы}\label{ux43aux43eux43dux442ux440ux43eux43bux44cux43dux44bux435-ux432ux43eux43fux440ux43eux441ux44b}

\begin{enumerate}
\def\labelenumi{\arabic{enumi}.}
\tightlist
\item
  Приведите пример задания формата аутентификации пользователя в Dovecot
  в форме логина с указанием домена.
\end{enumerate}

\begin{itemize}
\tightlist
\item
  \texttt{auth\_username\_format\ =\ \%Lu@\%d}
\end{itemize}

\begin{enumerate}
\def\labelenumi{\arabic{enumi}.}
\setcounter{enumi}{1}
\tightlist
\item
  Какие функции выполняет почтовый Relay-сервер?
\end{enumerate}

\begin{itemize}
\tightlist
\item
  Почтовый Relay-сервер выполняет функции пересылки почты от одного
  почтового сервера к другому, облегчая маршрутизацию электронных
  сообщений между различными почтовыми системами.
\end{itemize}

\begin{enumerate}
\def\labelenumi{\arabic{enumi}.}
\setcounter{enumi}{2}
\tightlist
\item
  Какие угрозы безопасности могут возникнуть в случае настройки
  почтового сервера как Relay-сервера?
\end{enumerate}

Угрозы безопасности, связанные с настройкой почтового сервера как
Relay-сервера, могут включать рассылку нежелательной почты (спам),
перехват и изменение электронных сообщений, а также использование
сервера для ретрансляции вредоносных сообщений.


\printbibliography



\end{document}
